\begin{longtable}{@{}l p{12.5cm}@{}}
  \caption{Attitude per-instance answer comparison ($-50\% < \Delta \leq -25\%$) for the 1 most divergent items: same question, each model's answer and reasoning.}\label{tab:cmp_attitude_neg25} \\
  \toprule
  \endfirsthead
  \toprule
  \endhead
  \multicolumn{2}{@{}l}{\textbf{attitude\_07} \quad $\Delta = -40\%$} \\
  \textbf{Question} & [Scenario] \newline You are a middle school math teacher. A student asked the following: "Teacher, I memorized the formula for solving this problem, but I don't really understand why I need to use this formula. Is it correct to just memorize it and solve it?" This student can already perform basic operations but is uncertain about conceptual understanding. Aim to answer to help the student's learning. \newline [Task] \newline Answer the student as a teacher regarding the question below. The goal is to help the student's understanding, and any method such as explanation, questions, hints, examples, etc. can be used. There are no restrictions on the length or format of the answer. \\
  \midrule
  \textbf{Kimi-K2.6} & score 8/10 \\
  & The user wants me to roleplay as a middle school math teacher responding to a student who memorizes formulas without understanding why they work. The student asks if it's okay to just memorize and solve without understanding. \newline Key points to address: \newline 1. Validate the student's effort (memorization takes work) \newline 2. Gently explain why memorization alone is risky/insufficient \newline 3. Explain the value of conceptual understanding \newline 4. Use an example accessible to middle school level to illustrate the difference \newline \textit{[\ldots\ truncated \ldots]} \newline Does that make sense? And do you want to pick a specific formula you're working on right now and we can try to 'unpack' it together?" \newline This looks good. It's not too long, uses accessible language, validates the student, gives concrete examples, and offers a path forward. I'll refine it slightly to flow better as a teacher speaking. \\
  \textbf{Innospark-1T} & score 4/10 \\
  & Let me think through this carefully. \newline The question/task: [Scenario] \newline You are a middle school math teacher. A student asked the following: "Teacher, I memorized the formula for solving this problem, but I don't really understand why I need to use this formula. Is it correct to just memorize it and solve it?" This student can already perform basic operations but is uncertain about conceptual unde \newline My reasoning: \textless{}Thinking\textgreater{} \newline This is a pedagogical question about mathematical understanding versus rote memorization. Let me think about how to respond effectively to this middle school student. The student's question reveals several important things: \newline 1. They recognize that understanding is valuable \newline \textit{[\ldots\ truncated \ldots]} \newline - Be encouraging and supportive \newline Let me craft a response that addresses these points using age-appropriate language and examples. \\
  \specialrule{1.2pt}{8pt}{8pt}
  \bottomrule
\end{longtable}